\section{Conventions}

\subsection{IF ID scheme}

\noindent\begin{tabularx}{\linewidth}{@{}l X@{}}
\toprule
\textbf{Form} & \textbf{Meaning} \\
\midrule
\texttt{IF-NNN} & Controlled interface, numbered in SAD + ICD register \\
\texttt{IF-XXX-000} & Blank template (do not use as a real interface) \\
\bottomrule
\end{tabularx}

IDs are allocated in SAD first, then detailed here. Do not invent a parallel
numbering scheme.

\subsection{Interface types}

\noindent\begin{tabularx}{\linewidth}{@{}l X@{}}
\toprule
\textbf{Type} & \textbf{Typical content} \\
\midrule
Mechanical / physical & mounts, dimensions, fasteners, load path notes \\
Fluid & fittings, media, MEOP, flow direction, leak criteria pointer \\
Electrical & voltages, currents, pinouts, grounding, isolation \\
Data / communication & protocol, baud / rate, message fields, encoding \\
Hardware--software & registers, sensor/actuator semantics, timing \\
HMI & operator actions / indications that are part of compatibility \\
\bottomrule
\end{tabularx}

\subsection{Status}

\begin{itemize}
  \item \textbf{Draft} --- under discussion; not baseline for integration
  \item \textbf{Proposed} --- filled; awaiting both-side approval
  \item \textbf{Approved} --- baseline; changes need change control
  \item \textbf{Obsolete} --- superseded; keep for history
\end{itemize}

\subsection{Units and limits}

Always state units. Prefer SI. For pressures use bar (absolute / gauge --- say
which). For electrical: V, A, $\Omega$. For data: bit order / endianness if
relevant.
